%%%%%%%%%%%%%%%%%%%%%%%%%%%%%%%%%%%%%%%
% Cover Letter - Waymo
% 2027 Summer Intern, BS/MS, Embedded, Software Engineer
% Plain XeLaTeX — matches simple resume style (Spectral font)
%%%%%%%%%%%%%%%%%%%%%%%%%%%%%%%%%%%%%%%

\documentclass[11pt]{article}

\usepackage[margin=0.75in, top=0.7in, bottom=0.7in]{geometry}
\usepackage[hidelinks]{hyperref}
\usepackage{parskip}
\usepackage{fontspec}

\setmainfont{Spectral}[
  Path = ../../cv/fonts/,
  UprightFont    = Spectral-Regular.ttf,
  BoldFont       = Spectral-Bold.ttf,
  ItalicFont     = Spectral-Italic.ttf,
  BoldItalicFont = Spectral-BoldItalic.ttf
]

\setlength{\parindent}{0pt}
\setlength{\parskip}{14pt}
\linespread{1.15}
\raggedright

% No hyphenation
\hyphenpenalty=10000
\exhyphenpenalty=10000

\pagestyle{empty}
\urlstyle{same}

\begin{document}

% ── DATE ─────────────────────────────────────────────────────────────────────
\begin{flushright}
  \today
\end{flushright}

% ── BODY ─────────────────────────────────────────────────────────────────────
Dear Waymo Hiring Team,

I'm applying for the internship on Waymo's Embedded Software Engineering
team for Summer 2027. I'm a Computer Science student at RIT in the
combined bachelor's and
master's program, with an expected graduation date of December 2027, and
have a strong interest in firmware and embedded software. At GE Vernova, I
built a full I/O subsystem from scratch in C on an STM32F7 for industrial
radios. On the RIT Electric Vehicle Team, I'm the vehicle control unit
(VCU) Software Lead. The firmware I've built so far is for our current
track motorcycle's VCU, bringing up a custom STM32F446 board and
integrating it with the rest of the vehicle over CAN and CANopen. In my
free time I'm building an autonomous fixed-wing flight controller on an
ESP32S3, with an 80-test regression suite covering step responses, noise
sweeps, and waypoint missions to catch edge cases before it ever flies.

What got me interested in Waymo specifically is how the compute redundancy
is built. Two independent compute units run the full workload in parallel,
and if one faults, the other takes over, backed by a secondary computer
whose only job is to bring the vehicle to a safe stop if something goes
wrong. That's the same direction my team is building toward on the VCU,
adding a second independent microcontroller that watches the main one and
can force a safe shutdown if it detects a problem. I want to help build
that kind of fault-tolerant embedded software.

Hardware bring-up and lab debugging are a big part of how I work. I debug
at the hardware level regularly with a Saleae logic analyzer, an
oscilloscope, and a J-Link Segger, and I've written peripheral drivers for
EVT's hardware-agnostic library directly from microcontroller datasheets
and reference manuals. My
flight controller project is the closest thing I have to the kind of
embedded work this role needs. I wrote the PID stabilization and navigation
firmware for the aircraft, and built SITL simulators alongside it to
validate the gains across 80 regression tests before testing on real
hardware. At Lockheed Martin, I wrote a
Python tool that automated QA checks across 100+ repositories and was
deployed as a mandatory GitLab CI gate, so I'm comfortable writing the kind
of test tooling and automation this role also needs.

I would welcome the opportunity to talk more about where I could
contribute. Thank you for your consideration.

Thanks,

Travis Brown

\end{document}
